\subsection{A characteristic-power tower has only finitely many eligible primes}
\label{sec:pro-two-tower-obstruction}

\paragraph{Finding A44 (a finite obstruction to the controlled-tower attempt).}
Let $A/\mathbf F_q$ be an abelian surface, with $q$ a power of two.
Write its Frobenius polynomial as
\[
 \begin{gathered}
 P_A(T)=T^4-aT^3+bT^2-qaT+q^2,\qquad a,b\in\mathbf Z,\\
 K=\bigcup_{m\ge0}\mathbf F_{q^{2^m}}.
 \end{gathered}
 \tag{CT.1}
\]
Call an odd prime $\ell$ eligible if
$\langle2,-1\rangle=(\mathbf Z/\ell)^\times$.
Define the following nonzero integers:
\[
 \begin{aligned}
 D_\epsilon&=(\epsilon b-q^2-1)(q+\epsilon)^2-qa^2,
                         &&\epsilon\in\{1,-1\},\\
 L_A&=\left|\#A(\mathbf F_{q^{16}})
                    (q^4-1)D_1D_{-1}\right|.
 \end{aligned}
 \tag{CT.2}
\]
Every eligible prime occurring in the order of a point of $A(K)$
divides $L_A$. Consequently there is a finite, computable set
\[
 S_A=\{\ell:\ell\text{ is odd},\ \ell\mid L_A,\
                 \langle2,-1\rangle=(\mathbf Z/\ell)^\times\}
 \tag{CT.3}
\]
containing every such prime. Put
\[
 e=\max_{\ell\in S_A}v_2\!\left(
                   \prod_{j=1}^{4}(\ell^j-1)\right),
 \tag{CT.4}
\]
with $e=0$ if $S_A$ is empty. Then
\[
          A(K)_{S_A}=A(\mathbf F_{q^{2^e}})_{S_A}.
 \tag{CT.5}
\]
In particular, all points of $A(K)$ whose odd order $n>1$ satisfies
$\langle2,-1\rangle=(\mathbf Z/n)^\times$ belong to the finite
group on the right of (CT.5). This is an unconditional statement.

For a genus-two curve in characteristic two, the actual Raynaud
theta support nevertheless contains infinitely many nonzero
odd-order points over this same field $K$. All but finitely many
of them therefore have a prime factor in their order which is
individually ineligible. Thus abundance in this constant-field
tower cannot by itself supply an unbounded family of the orders
required by the single-orbit primitive bridge.

\paragraph{Known inputs and exact scope.}
We use the Frobenius polynomial, its point-count formula and its
Weil absolute values \cite[Proposition~12.9 and Theorem~19.1]{MilneAV}, and the
prime-to-characteristic torsion module
\cite[Theorem~8.2 and Remark~8.4]{MilneAV}.
The finite-primary stability in a tower is consistent with
Poonen's Lemma~6.2 \cite{Poonen2005}; the new calculation here
isolates a finite set of eligible primes in the characteristic-two
tower. The argument does not treat towers of odd auxiliary-prime
degree. Nor does finiteness say that the remaining finite set of
theta witnesses is empty: the example below is nonempty.

\begin{proof}
\emph{1. The integers in (CT.2) do not vanish.}
All complex roots of $P_A$ have absolute value $\sqrt q>1$.
If $D_\epsilon=0$, put $s=a/(1+\epsilon q)$. Direct coefficient
comparison would give a factorization over $\mathbf Q$,
\[
 P_A(T)=(T^2-sT+\epsilon)
                 (T^2-\epsilon qsT+\epsilon q^2).
 \tag{CT.6}
\]
The two roots of the first factor would be roots of $P_A$,
so their product would have absolute value $q$. Its constant
term instead has absolute value one. This is a contradiction.
The other factors in $L_A$ are positive nonzero integers.

\emph{2. A prime first appearing late supplies a cyclotomic eigenvalue.}
Let an eligible prime $\ell$ divide $\#A(\mathbf F_{q^{2^m}})$,
with $m$ least. If $m\le4$, then
$\ell\mid\#A(\mathbf F_{q^{16}})$ and it is already in (CT.2).
Suppose $m\ge5$. On the four-dimensional space
$A[\ell](\overline{\mathbf F}_q)$, Frobenius $\Phi$ has an
eigenvalue $\lambda$ of exact multiplicative order $2^m$.
Indeed the nonzero kernel of $\Phi^{2^m}-1$ is semisimple for
$\Phi$, since $T^{2^m}-1$ is separable modulo $\ell$.
If all eigenvalues in that kernel had order dividing $2^{m-1}$,
there would be nonzero rational $\ell$-torsion at the preceding
level, contrary to minimality. The minimal polynomial of
$\lambda$ over $\mathbf F_\ell$ divides $P_A\bmod\ell$; hence
\[
       [\mathbf F_\ell(\lambda):\mathbf F_\ell]
               =\operatorname{ord}_{2^m}(\ell)\le4.
 \tag{CT.7}
\]
This degree is a power of two, so it is $1$, $2$ or $4$.
No semisimplicity of the full Frobenius matrix is assumed.

Eligibility excludes $\ell\equiv1\pmod8$: in that residue class
both $2$ and $-1$ are squares modulo $\ell$, so they cannot
generate the full unit group. Here the assertion for $2$ follows
directly from Gauss's lemma: for $\ell=8u+1$ the number of
$j$ with $1\le j\le(\ell-1)/2$ and $2j>\ell/2$ is $2u$.
If $\ell\equiv3$ or $5\pmod8$, then
$v_2(\ell^4-1)=4$, contradicting (CT.7) when $m\ge5$.
Consequently
\[
                              \ell\equiv7\pmod8.
 \tag{CT.8}
\]
The degree in (CT.7) cannot be one, because then
$\ell\equiv1\pmod{2^m}$.

\emph{3. Degree two forces one of two fixed divisibilities.}
Suppose $\lambda$ has degree two, and write
$t=v_2(\ell+1)$; by (CT.8), $t\ge3$ and $v_2(\ell-1)=1$.
The condition $2^m\mid\ell^2-1$ implies $t\ge m-1$.
Therefore
\[
          \lambda^{\ell+1}=\epsilon\in\{1,-1\},
       \qquad
          \operatorname{min}_\lambda(T)=T^2-sT+\epsilon.
 \tag{CT.9}
\]
The Frobenius polynomial is $q$-reciprocal, so its roots also
include $q/\lambda$ and $q/\lambda^\ell$. If these two roots
overlap the first pair, then either $q=\epsilon$ in
$\mathbf F_\ell$, or $q=\lambda^2$.
In the latter case $\lambda^2$ is in $\mathbf F_\ell$ and
$\lambda\notin\mathbf F_\ell$, so
$\lambda^\ell=-\lambda$ and $\epsilon=-q$.
Thus every overlap case gives $\ell\mid q^2-1$, already a factor
of $L_A$.

In the remaining case all four roots are distinct and exhaust
$P_A\bmod\ell$. The factorization (CT.6) holds modulo $\ell$.
Its coefficients give
\[
        a=(1+\epsilon q)s,\qquad
        b=\epsilon(1+q^2+qs^2).
 \tag{CT.10}
\]
Eliminating $s$ without division gives $\ell\mid D_\epsilon$.
This proves the required assertion in degree two.

\emph{4. Degree four forces another fixed divisibility.}
If $\lambda$ has degree four, then
$2^m\mid\ell^4-1$ but $2^m\nmid\ell^2-1$.
Since $v_2(\ell^2+1)=1$, this forces
$v_2(\ell+1)=m-2$. Hence
\[
 v_2(1+\ell+\ell^2+\ell^3)
       =v_2((1+\ell)(1+\ell^2))=m-1.
\]
The norm of $\lambda$ to $\mathbf F_\ell$ is therefore $-1$.
Its degree-four minimal polynomial is $P_A\bmod\ell$, whose
constant term is $q^2$. We conclude that
$\ell\mid q^2+1$, again a factor of $L_A$.
Together with Steps~2--3 this proves (CT.3).

\emph{5. Arbitrary primary exponents are already in one finite field.}
Let $\ell\in S_A$ and let $x\in A(K)$ have order $\ell^r$.
Its field degree over $\mathbf F_q$ is a power of two, and is
the length of its Frobenius orbit. This divides the order of
the Frobenius matrix in
$\operatorname{GL}_4(\mathbf Z/\ell^r)$ and hence the order of
that finite group. Reduction modulo $\ell$ has an $\ell$-group
kernel. As $\ell$ is odd, the two-part of its order is exactly
that of $\operatorname{GL}_4(\mathbf F_\ell)$, namely
\[
        2^{\,v_2(\prod_{j=1}^4(\ell^j-1))}.
 \tag{CT.11}
\]
The orbit degree thus divides $2^e$. Adding the finitely many
primary components proves (CT.5). Reduction of the unit group
modulo any prime divisor of an odd $n$ is surjective; therefore
the full signed-unit condition at $n$ implies eligibility at
every such prime. This proves the finite-field assertion for
all the stated orders, with no bound on their exponents.

\emph{6. Odd theta points in the tower are still unbounded.}
We prove this last statement so that the obstruction does not
merely come from a lack of odd torsion. First, for any abelian
surface as in (CT.1), put
\[
 d_m=\#A(\mathbf F_{q^{2^m}})[2^\infty].
\]
There are $m_0\ge0$ and an integer $0\le t\le2$ such that
\[
                     d_m=d_{m_0}2^{t(m-m_0)}
                         \qquad(m\ge m_0).
 \tag{CT.12}
\]
To see this, put all four roots $\alpha_i$ of $P_A$ in a finite
extension of $\mathbf Q_2$, with valuation normalized by $v(2)=1$.
The point-count product gives
$v_2(d_m)=\sum_i v(1-\alpha_i^{2^m})$.
The roots are integral. A nonunit contributes zero. A unit whose
residue is not one also contributes zero: its residue has odd
order, which cannot divide $2^m$.
For a unit with residue one, repeated squaring eventually gives
$v(\alpha_i^{2^{m_0}}-1)>1$. Explicitly, while the positive
valuation is less than one, squaring doubles it; at valuation
one the next valuation is at least two. Thereafter
\[
  v(u^2-1)=v(u-1)+1\quad\text{when }v(u-1)>1.
\]
No root is a root of unity, by its complex absolute value.
Thus each such unit contributes an affine function of $m$ with
slope one for all sufficiently large $m$.
The number of unit roots is at most two: the reciprocal pairing
$\alpha\mapsto q/\alpha$ sends them injectively to nonunit
roots, with multiplicities. Summing proves (CT.12).

Now take $A=J_X$, where $X=C^{(1)}$ and
$\Theta=i_B(X)$ is the actual genus-two Raynaud theta support.
At the individual field of size $Q_m=q^{2^m}$, use the Lang
quotient whose deck group is the rational two-primary subgroup.
The argument of A12, or Steps~2--5 of A42 at this single field,
gives
\[
 \begin{aligned}
 O_m&=\#\bigl(\Theta(\mathbf F_{Q_m})
                  \cap J_X(\mathbf F_{Q_m})_{\rm odd}\bigr),\\
 O_m&\ge\frac{Q_m+1}{d_m}
                      -2\left(1+\frac1{d_m}\right)\sqrt{Q_m}.
 \end{aligned}
 \tag{CT.13}
\]
Here the counting covers may vary with $m$; A42's fixed-primary
extension hypothesis is not being asserted. The bound (CT.12)
is polynomial in $2^m$, whereas $Q_m$ is exponential in $2^m$.
The right side of (CT.13) tends to infinity. There is at most
one zero theta class, so this gives infinitely many nonzero
odd theta points over $K$.
The group of points of $J_X(K)$ supported on eligible primes
is finite by (CT.3)--(CT.5). Every other odd theta point must
therefore have an ineligible prime factor in its order.
\end{proof}

\paragraph{A finite candidate set can contain useful points.}
For the A13 curve
\[
 y^2+y=x+x^{-1}+(x+1)^{-1}
\]
over $\mathbf F_2$, $P_A(T)=T^4-T^2+4$, so $a=0$, $b=-1$,
$D_1=-54$ and $D_{-1}=-4$. The established quadratic recurrence
for $4$-power Frobenius gives
\[
 \#A(\mathbf F_{2^{16}})=65760^2
                    =2^{10}3^2 5^2 137^2.
 \tag{CT.14}
\]
Since $137\equiv1\pmod8$, (CT.3) gives $S_A=\{3,5\}$.
In particular the order-$4159$ witness of A13 lies outside this
constant-field tower; it was constructed over an odd-degree
extension.

There are in fact order-three theta points in the tower.
Put $h=X^2+X$ and use the polynomial model
\[
 Y^2+hY=X^5+X^3+X^2+X=:f.
\]
Let $W_0,W_1,W_\infty$ be its three Weierstrass points.
The curve and its Frobenius twist have this same
$\mathbf F_2$-model. The Raynaud theta characteristic is
$B=W_0+W_1-W_\infty$.
Here are the details of this identification. The three simple
Artin--Schreier poles make the genus-two curve ordinary, so
$J[2](\overline{\mathbf F}_2)$ has four elements and there are
exactly four geometric theta characteristics. The three $W_i$
are distinct effective ones. Since $2W_i\sim K_C$, the displayed
$B$ is another theta characteristic. It differs from $W_0$
and $W_1$ because distinct degree-one point divisors are not
linearly equivalent. If $B\sim W_\infty$, then
$W_0+W_1\sim2W_\infty$ would be a canonical hyperelliptic fibre,
impossible for two distinct branch points. Thus $B$ is the
fourth characteristic. Any effective characteristic is a point
$R$ with $2R\sim K_C$, hence a Weierstrass point. Consequently
$B$ is ineffective. The locally exact differential
characteristic is ineffective on an ordinary curve, by the
Frobenius cohomology sequence, so it is precisely this $B$.
All three $W_i$ are rational over $\mathbf F_2$.

Choose $r\in\mathbf F_4$ with $r^2+r+1=0$, and
$a_0\in\mathbf F_{16}$ with $a_0^2+a_0=r^2$.
Then $R=(r,a_0)$ lies on the curve, and direct multiplication gives
\[
 \begin{aligned}
 F&=Y+a_0h,\\
 F\,\iota(F)&=f+(a_0^2+a_0)h^2
                          =h(X+r)^3,\\
 \operatorname{div}(F)&=3R+W_0+W_1-5W_\infty.
 \end{aligned}
 \tag{CT.15}
\]
The conjugate value at $R$ is one, since $h(r)=1$; the pole of
$F$ at infinity has order five. These facts and the norm identity
justify the full divisor formula. Since
$\operatorname{div}(h)=2W_0+2W_1-4W_\infty$, it follows that
\[
                    \operatorname{div}(F/h^2)=3(R-B).
 \tag{CT.16}
\]
The class $R-B$ is nonzero because $B$ is ineffective.
This gives four exact order-three theta points over
$\mathbf F_{16}$, one for each of the two choices of $r$ and
two choices of $a_0$. Thus the finite set in A44 need not be
empty, and a finite member can already suffice for A11.
The proof does not assert that these are all eligible theta
points of this curve.

\paragraph{Status of the attempted main-route application.}
The computation closes the specific proposed inference from
unbounded odd theta torsion in a characteristic-power constant
tower to an unbounded supply of signed-unit-compatible orders.
It leaves a finite, decidable candidate set for that tower,
which can contain successful witnesses as (CT.15) shows.
It does not settle the existence of a witness in that set for
every curve, and it does not rule out an odd auxiliary-prime
tower, a different cover, or another primitive-factor mechanism.
No previously open universal prescribed-point row is marked
closed by this obstruction.
